\documentclass[10pt, a4paper]{aqademic}

\usepackage[spanish]{babel}
	\selectlanguage{spanish}

% Document packages

\usepackage{amsmath}
\usepackage{amsfonts}
\usepackage[type=CC, modifier=by-nc-sa, version=4.0]{doclicense}
\usepackage{graphicx}
\usepackage{tikz}

% Document settings

\author{Atanasio José Rubio Gil}
\title{Álgebra lineal y estructuras matemáticas}

\AqSetChapter{Tema}

\newcommand\Division[2]{%
	$\strut#1$\kern.25em\smash{\raise-.35ex\hbox{\rule{0.4pt}{2ex}}}$\mkern-0.72mu
	\underline{\,#2}$
}

\newcount\total
\newcount\lasttotal
\newcount\targetbase

% https://tex.stackexchange.com/questions/107637/repeated-division-converting-from-base-10-to-another-base
\def\CambioDeBase#1#2{%
	\begin{tikzpicture}[every node/.style={minimum width=1cm, minimum height=0.5cm}, x=1cm,y=0.5cm]
	%
	\total=#1%
	\targetbase=#2
	\def\newnumber{}
	%
	\pgfmathloop
	\ifnum\total<1
	\else
		%
		\ifnum\pgfmathcounter>1
			\node at (\pgfmathcounter, -\pgfmathcounter+1) (tmp) {\the\targetbase};
			\draw (tmp.north west) |- (tmp.south east);
			%
			\node at (\pgfmathcounter-1, -\pgfmathcounter) (tmp) {\pgfmathparse{int(\total*\targetbase)}\pgfmathresult};
			\draw (tmp.south west) -- (tmp.south east);
			%
			\pgfmathparse{int(\lasttotal-\total*\targetbase)}%
			\let\digit=\pgfmathresult
			\node at (\pgfmathcounter-1, -\pgfmathcounter-1) [text=red] {\digit};
			\edef\newnumber{\digit\newnumber}
		\fi
		%
		\ifnum\total<\targetbase
			\edef\newnumber{\the\total\newnumber}
			\node at (\pgfmathcounter, -\pgfmathcounter) [text=red]  {\the\total};
		\else
			\node at (\pgfmathcounter, -\pgfmathcounter) {\the\total};
		\fi
		\lasttotal=\total
		\divide\total by\targetbase
	\repeatpgfmathloop
	\node [anchor=west] at (1, -\pgfmathcounter-1) {$#1=\newnumber_{\the\targetbase}$};
	\end{tikzpicture}
}

% Document composition

\begin{document}

\AqMaketitle[%
	cover    = logo-ugr.png,
	org      = Grado en Ingeniería Informática,
	subtitle = Teoría,
	url      = https://github.com/Groctel/ugr-informatica
]

\begin{titlepage}

\newgeometry{%
	left=5cm,
	right=5cm
}

\thispagestyle{empty}

\topskip0pt
\vspace*{\fill}

\begin{center}
	\textbf{\Huge{¡Gracias por apostar por el conocimiento libre!}}
\end{center}

\vspace{1cm}

Estos apuntes están compartidos libremente para que puedas estudiar de forma gratuita y sin bloques de publicidad molesta en tus páginas.
Su publicación y uso son un paso adelante en la idea de que el conocimiento debería ser compartido de forma libre y gratuita para todos sin ningún tipo de impedimentos o distracciones generadas por muros de pago.

Como autor, no recibo ningún tipo de ingreso por la confección y publicación de este material.
Si quieres colaborar y ayudarme a seguir ofreciéndote facilidades para estudiar, puedes unirte al proyecto a través del enlace de la portada o enviarme una propina mediante el siguiente botón o su enlace:

\vspace{0.5cm}

\begin{center}
	\href{https://ko-fi.com/groctel}{\includegraphics[scale=0.15]{RecursosTeX/BuyMeACoffee_blue-2x.pdf}}
	\url{https://ko-fi.com/groctel}
\end{center}

\vspace{0.5cm}

Dado que este documento está compartido bajo una licencia \textbf{CC-by-nc-sa}, tienes la libertad de distribuirlo y adaptarlo bajo las siguientes condiciones:

\begin{itemize}
	\item\textbf{BY:} Debes darme crédito de forma adecuada con un enlace a la licencia e indicar si has hecho cambios.
	\item\textbf{NC:} No puedes hacer uso del mismo con propósitos comerciales.
	\item\textbf{SA:} Debes distribuir tus modificaciones bajo la misma licencia que el documento original.
\end{itemize}

Puedes leer la licencia completa pinchando en el enlace de la portada.

\vspace*{\fill}

\end{titlepage}

\tableofcontents

\chapter{Números naturales y números enteros}
	\section{Representación de los números naturales}\label{representacion-numeros-naturales}

\subsection{$\mathbb{N}$, el conjunto de los números naturales}\label{el-conjunto-de-los-numeros-naturales}

Los números naturales es un conjunto infinito de números que, coloquialmente, podemos enumerar como $0,1,2,\ldots$ hasta aburrirnos.
Si queremos hacerlo matemáticamente, debemos describirlo en función de las propiedades que satisface.
Enumerémoslas:

\subsubsection{Los números naturales se pueden sumar}

Dados dos números naturales $m$ y $n$, hay definido un nuevo número natural que es la suma de $m$ y $n$ y se representa como $m+n$.
Hacemos hincapié en que el nuevo número es natural o, expresado matemáticamente:

\[m,n \in\mathbb{N} \Rightarrow m+n \in\mathbb{N}, \forall m,n\]

Esto es importante porque la inversa de la suma, la resta, no siempre cumple esta propiedad.
Por ejemplo, $4-5=-1$ y éste no es un número natural.
La suma satisface las siguientes propiedades:

\begin{center}
\begin{tabular}{l l}
	\textbf{Propiedad}  & \textbf{Expresión}                                       \\
	\toprule
	Asociativa          & $(m+n) + p = m + (n+p), \forall m,n,p \in\mathbb{N}$     \\
	Conmutativa         & $m+n = n+m, \forall m,n \in\mathbb{N}$                   \\
	Elemento neutro (0) & $m+0 = m, \forall m \in\mathbb{N}$                       \\
	Cancelativa         & $m+n = m+p \Rightarrow n=p, \forall m,n,p \in\mathbb{N}$ \\
\end{tabular}
\end{center}

\subsubsection{Los números naturales se pueden multiplicar}

Dados dos números naturales $m$ y $n$, hay definido un nuevo número natural que es el producto de $m$ y $n$ y se representa como $m \cdot n$ o simplemente $mn$.
De nuevo, lo expresamos matemáticamente de la siguiente forma:

\[m,n \in\mathbb{N} \Rightarrow m \cdot n \in\mathbb{N}, \forall m,n\]

Al igual que la resta no cumple siempre esta propiedad, la división tampoco lo hace.
Por ejemplo, $\frac{4}{5}=0.8$ y éste no es un número natural.
La multiplicación satisface las siguientes propiedades:

\begin{center}
\begin{tabular}{l l}
	\textbf{Propiedad}  & \textbf{Expresión}                                                             \\
	\toprule
	Asociativa          & $(m \cdot n) \cdot p = m \cdot (n \cdot p), \forall m,n,p \in\mathbb{N}$       \\
	Conmutativa         & $m \cdot n = n \cdot m, \forall m,n \in\mathbb{N}$                             \\
	Elemento neutro (1) & $m \cdot 1 = m, \forall m \in\mathbb{N}$                                       \\
	Cancelativa         & $m \cdot n = m \cdot p \Rightarrow n=p, \forall m,n,p \in\mathbb{N}, m \neq 0$ \\
\end{tabular}
\end{center}

También tenemos que la suma es distributiva respecto al producto:

\[m \cdot (n+p) = m \cdot n + m \cdot p, \forall m,n,p \in\mathbb{N}\]

\subsubsection{Los números naturales se pueden ordenar}

Cuando sumamos dos números naturales obtenemos un número mayor a los dos números anteriores.
Expresamos esta relación de la siguiente forma:

\[\exists p \in\mathbb{N} : m+p = n \Rightarrow m \leq n, \forall m,n \in\mathbb{N}\]

Esta relación de orden (\textit{menor o igual que} o $\leq$) satisface las siguientes tres propiedades con nombre\footnote{%
	En general cualquier conjunto con cualquier relación de orden cumple estas tres propiedades, aunque aquí nos estemos centrando en el conjunto $\mathbb{N}$ y la relación $\leq$.
}:

\begin{center}
\begin{tabular}{l l}
	\textbf{Propiedad} & \textbf{Expresión}                                                          \\
	\toprule
	Reflexiva          & $m \leq m, \forall m \in\mathbb{N}$                                         \\
	Asimétrica         & $m \leq n \land n \leq m \Rightarrow m=n, \forall m,n \in\mathbb{N}$        \\
	Transitiva         & $m \leq n \land n \leq p \Rightarrow m \leq p, \forall m,n,p \in\mathbb{N}$ \\
	Orden total        & $m \leq n \lor n \leq m, \forall m,n \in\mathbb{N}$                         \\
\end{tabular}
\end{center}

Esta relación también cumple las siguientes propiedades

\begin{itemize}
	\item $m \leq n \Rightarrow m+p \leq n+p, \forall m,n,p \in\mathbb{N}$
	\item $m+p \leq n+p \Rightarrow m \leq n, \forall m,n,p \in\mathbb{N}$
	\item $m \leq n \Rightarrow m \cdot p \leq n \cdot p, \forall m,n,p \in\mathbb{N}$
	\item $m \cdot p \leq n \cdot p \Rightarrow m \leq p, \forall m,n,p \in\mathbb{N}$
\end{itemize}

Sin embargo, todas las propiedades enumeradas hasta ahora se cumplen para otros conjuntos, como $\mathbb{Q}^+$ o $\mathbb{R}^+$.
Para definir los números naturales necesitamos una propiedad que no se cumpla en ningún otro conjunto.

\subsection{El principio de induccion}\label{el-principio-de-induccion}

\[\{A \subseteq \mathbb{N} : 0 \in A, n \in A \Rightarrow n+1 \in A\} \Rightarrow A = \mathbb{N}\]

Desglosemos este principio paso por paso.
Primero, tomemos un conjunto $A \subseteq \mathbb{N}$, es decir, que $A$ sea un subconjunto de $\mathbb{N}$ pudiendo ser todo $\mathbb{N}$.
Segundo, tenemos que $0 \in A$.
Como el 0 es el menor de los números naturales\footnote{%
	En algunas disciplinas se considera que el menor número natural es el 1 porque el 0 les causa demasiado inconveniente y es más fácil obviarlo que explicitar constantemente que lo escrito se aplica a $\mathbb{N}$\textbackslash0.
}, hacemos que $A$ lo contenga y lo preparamos para el truco final:
La existencia de un número $n$ implica necesariamente la existencia de $n+1$.
Como nuestro conjunto contiene el 0, debe contener necesariamente el 1.
Como ahora contiene el 1, debe contener necesariamente el 2, que da paso al 3, al 4, al 5\ldots
Si nuestro conjunto contiene el 52543543, también contendrá necesariamente el 52543544.
Siguiendo esta progresión, podemos contar un número infinito de números, por lo que $A$ contiene todos los números naturales y se verifica que $A=\mathbb{N}$.

Este proceso puede compararse análogamente al efecto de derribar piezas de dominó colocadas secuencialmente de pie.
Para que caiga la pieza $p_n$ no buscamos que caiga la pieza $p_{n-1}$ y ésta derribe a la pieza $p_n$, nos basta con derribar manualmente la pieza $p_0$ ($0 \in A$) y cada pieza $p_n$ derribe a la pieza $p_{n+1}$.

Una consecuencia de este principio es el \textbf{Principio de buena ordenación}:

\[A \subseteq \mathbb{N} : A \neq \emptyset \Rightarrow A\text{ tiene mínimo}\]

Como corolario a este principio, tenemos que no existen sucesiones infinitas y estrictamente decrecientes en $\mathbb{N}$, ya que en algún momento deberá llegar necesariamente al mínimo.

\pagebreak

\subsection{Otras propiedades de los números naturales}

\subsubsection{Los números naturales pueden elevarse}

De la misma forma que la multiplicación es la aplicación repetida de la suma, tenemos que la potencia es la aplicación repetida de la multiplicación, de forma que, dados dos números naturales $m$ y $n$, hay definido un nuevo número natural que es la potencia de $m$ y $n$ (o $m$ elevado a $n$) y se representa como $m^n$ siempre y cuando no se cumpla que $m=n=0$.
Expresamos matemáticamente esta operación de la siguiente forma:

\[m,n \in\mathbb{N} \Rightarrow m^n \in\mathbb{N}, \forall m,n \neq 0\]

Esta operación satisface las siguientes propiedades:

\begin{itemize}
	\item $m^{n+p} = nm^n \cdot m^p, \forall m,n,p \in\mathbb{N}, m \neq 0$
	\item ${(m^n)}^p = m^{n \cdot p}, \forall m,n,p \in\mathbb{N}, m,n \neq 0$
	\item ${(m \cdot n)}^p = mP p \cdot n^p, \forall m,n,p \in\mathbb{N}, m,n \neq 0$
\end{itemize}

\subsubsection{Existe un algoritmo para dividir los números naturales}

\[a,b \in\mathbb{N} \Rightarrow \exists! c,r \in\mathbb{N} : a = bc + r, 0 \leq r < b, b \neq 0\]

Este algoritmo es la \textit{división con caja} que conocemos desde primaria.
Como ejemplo, tengamos un dividendo $a=164$ y un divisor $b=7$:

\setlength{\parindent}{8cm}
\parbox{10cm}{%
	\Division{164}{7} \\
	024\hspace{1.05ex}23 \\
	\hspace*{1.1ex}\underline{03} \\
}
\setlength{\parindent}{15pt}

El algoritmo nos devuelve un cociente $c=23$ y un resto $r=3$.
Comparándolo con la expresión anterior, $a = bc + r \Leftrightarrow 164 = 7 \cdot 23 + 3$.
Con este método dividir números naturales teniendo siempre en cuenta el resto de la división.
Es muy importante tener en cuenta que la división entre 0 no está definida en $\mathbb{N}$. Otra forma de explicitar los números $c$ y $r$ es mediante las expresiones $(a\ div\ b)$ y $(a\ m\acute{o}d\ b)$ (módulo) respectivamente.

\subsection{Representación de los números naturales}

Al igual que tenemos diez dedos en las manos, los números que escribimos tienen diez cifras, del 0 al 9\footnote{%
	En algún sitio leí que los griegos se plantearon contar en base duodecimal porque tenemos doce falanges en los dedos, pero vieron que reeducar a todo el mundo era demasiado complicado a pesar de las ventajas del nuevo sistema.\textsuperscript{[\href{https://en.wikipedia.org/wiki/Wikipedia:Citation_needed}{Cita requerida}]}
}.
Es lo que llamamos base 10 o decimal.
La base nos indica el valor que puede adquirir cada cifra del número y el valor que se le otorga en función de su posición.
Expresémoslo formalmente de forma general:

\begin{align*}
	& b \in\mathbb{N} : b \geq 2 \\
	& a \in\mathbb{N} : a \neq 0 \\
	& \exists! k, a_0, a_1,\ldots,a_n \in\mathbb{N} : 0 \leq a_i < b \\
	& a = \sum_{k=0}^{m} a_kb^k = a_mb^m + a_{m-1}b^{k-1} + \cdots + a_1b + a_0 \\
	& a = a_k, a_{k-1}, \ldots, a_1, a_0 \big)_b
\end{align*}

Vamos paso por paso.
Primero, nombremos una base $b \geq 2$.
No podemos usar la base 1 porque, como veremos tras el segundo paso, sólo podríamos usar el 0 para escribir el número.
Segundo, tomemos un número $a \neq 0$, pues el 0 siempre se expresa como tal.
Con estas dos presunciones, tenemos que existe una única expresión de un número $a_0, a_1,\ldots,a_n$ de $k$ cifras, ambos naturales y de forma que cada cifra sea menor que la base, que cumpla que el número $a$ es la suma de todas sus cifras en función de su posición.
El valor que toma cada cifra equivale al valor intrínseco de la misma multiplicado por la base elevada a la posición en la que se encuentra la cifra de izquierda a derecha, siendo la primera 0.
Por ejemplo, el número 123 en base 10 tiene el valor de tres unos ($10^0$) más dos dieces ($10^1$) más un cien ($10^2$).
Ésta es la única forma de representar el 123 en base decimal, al igual que hay una única forma de representar un valor en cada base.

Cuando expresamos un valor en una base distinta a 10 (aunque podríamos hacerlo para base 10, pero sería demasiado esfuerzo) explicitamos la base con un paréntesis seguido de la base en subíndice.
Por ejemplo, para indicar que estamos escribiendo el número 174 en base 8 lo escribimos como $174\big)_8$.

\subsubsection{Cambios de base}

Para realizar el cambio de una base cualquiera a base decimal nos basta con utilizar la expresión anterior y calcular el valor decimal que le damos al número:

\[5243\big)6 = 5 \cdot 6^3 + 2 \cdot 6^2 + 4 \cdot 6 + 3 = 1179\big)_{10}\]

Para realizar el cambio de base decimal a una base cualquiera dividimos recursivamente hasta obtener un resto menor que la base y escribimos los restos de derecha a izquierda en orden de aparición:

\begin{center}
	\CambioDeBase{543}{4}
\end{center}

¿Qué ocurre si queremos hacer un cambio de un número $a$ en una base no decimal $b_1$ a un número $b$ en otra base no decimal $b_2$?
La opción difícil sería aprenderse las tablas de multiplicar en la base con la que operemos.
La opción fácil es tomar el camino largo y pasar primero el número $a$ a un tercer número $c$ base 10 y de ahí a la nueva base:

\begin{center}
\begin{tikzpicture}[scale=0.2]
\tikzstyle{every node}+=[inner sep=0pt]
	\draw (28.4,-22.1) node {$a\big)_{b_1}$};
	\draw (40.8,-22.1) node {$b\big)_{b_2}$};
	\draw (34.6,-30.5) node {$c\big)_{10}$};
	\draw [black] (30.18,-24.51) -- (32.82,-28.09);
	\fill [black] (32.82,-28.09) -- (32.75,-27.15) -- (31.94,-27.74);
	\draw [black] (36.38,-28.09) -- (39.02,-24.51);
	\fill [black] (39.02,-24.51) -- (38.14,-24.86) -- (38.95,-25.45);
	\draw [black] [dashed] (31.4,-22.1) -- (37.8,-22.1);
	\fill [black] (37.8,-22.1) -- (37,-21.6) -- (37,-22.6);
\end{tikzpicture}
\end{center}

Por otro lado, si queremos cambiar de un número en base $b$ a otro número en base $b^k$, nos basta con agrupar las cifras del número de derecha a izquierda de $k$ en $k$ y expresar dichos números en la nueva base:

\[01111010010\big)_2 = 0011 1101 0010\big)_2 = 3D2\big)_{16}\]


\end{document}
